\documentclass[11pt]{article}
\usepackage{mathblatt}
% gesetzt von werkzeuge/setzer.py; Lösungen nach bau/bauregeln.md „Lösungen“.
\newcommand{\kennung}{A5D}
\begin{document}
\pfheftstil
\fancyfoot[C]{{\footnotesize\color{mbgrau}\kennung\hfill\thepage}}
\pfheftkopf{Kürzen und Erweitern}{Klasse 6 \textperiodcentered{} Lösungen}

\begin{pfloesung}{}
\lz{1b)}{$\frac{4}{10}$ (4 Teile färben)}{$\frac{4}{10}$: Jedes Fünftel ist unten in 2 Teile geteilt, also $2 \cdot 2 = 4$ Teile färben.}{}
\end{pfloesung}
\begin{pfloesung}{}
\lz{2.}{a) mal 5 \quad b) durch 7 \quad c) mal 8}{a) erweitert mit $45 : 9 = 5$; b) gekürzt mit $28 : 4 = 7$; c) erweitert mit $64 : 8 = 8$}{}
\end{pfloesung}
\begin{pfloesung}{}
\lz{3.}{a) $\frac{12}{20}$ \quad b) $\frac{15}{18}$ \quad c) $\frac{70}{90}$}{a) $\frac{3 \cdot 4}{5 \cdot 4} = \frac{12}{20}$; b) $\frac{5 \cdot 3}{6 \cdot 3} = \frac{15}{18}$; c) $\frac{7 \cdot 10}{9 \cdot 10} = \frac{70}{90}$}{}
\end{pfloesung}
\begin{pfloesung}{}
\lz{4.}{a) $\frac{5}{6}$ \quad b) $\frac{3}{7}$ \quad c) $\frac{4}{7}$}{a) $\frac{20 : 4}{24 : 4} = \frac{5}{6}$; b) $\frac{15 : 5}{35 : 5} = \frac{3}{7}$; c) $\frac{24 : 6}{42 : 6} = \frac{4}{7}$}{}
\end{pfloesung}
\begin{pfloesung}{}
\lz{5.}{gleich viel}{Beide gleich viel: Lena $\frac{3}{4}$, Tom $\frac{9}{12} = \frac{9 : 3}{12 : 3} = \frac{3}{4}$.}{}
\end{pfloesung}
\begin{pfloesung}{}
\lz{6b)}{a) $\frac{3}{4}$ \quad b) $\frac{2}{3}$ \quad c) $\frac{2}{5}$}{a) $\frac{6}{8} = \frac{3}{4}$ (durch 2); b) $\frac{10}{15} = \frac{2}{3}$ (durch 5); c) $\frac{8}{20} = \frac{2}{5}$ (durch 4)}{}
\end{pfloesung}
\begin{pfloesung}{}
\lz{7.}{a) $\frac{3}{4}$ \quad b) $\frac{3}{5}$ \quad c) $\frac{2}{3}$ \quad d) $\frac{3}{4}$}{a) $\frac{18}{24} = \frac{9}{12} = \frac{3}{4}$ (durch 6); b) $\frac{36}{60} = \frac{18}{30} = \frac{3}{5}$ (durch 12); c) $\frac{48}{72} = \frac{24}{36} = \frac{2}{3}$ (durch 24); d) größte Zahl, die 36 und 48 teilt: 12; $\frac{36}{48} = \frac{36 : 12}{48 : 12} = \frac{3}{4}$}{}
\end{pfloesung}
\begin{pfloesung}{}
\lz{8.}{$\frac{16}{24} = \frac{2}{3}$}{16 von 24 Kästchen sind grau; durch 8 gekürzt.}{}
\end{pfloesung}
\begin{pfloesung}{}
\lz{9.}{a) $\frac{5}{6}$ \quad b) $\frac{3}{4}$ \quad c) $\frac{7}{20}$}{a) 1 h = 60 min: $\frac{50}{60} = \frac{5}{6}$; b) 1 kg = 1\,000 g: $\frac{750}{1\,000} = \frac{3}{4}$; c) 1 m = 100 cm: $\frac{35}{100} = \frac{7}{20}$}{}
\end{pfloesung}
\begin{pfloesung}{}
\lz{10.}{Ja}{$\frac{21}{28} = \frac{21 : 7}{28 : 7} = \frac{3}{4}$.}{}
\end{pfloesung}
\begin{pfloesung}{}
\lz{11b)}{a) $\frac{16}{24}$ \quad b) $\frac{15}{24}$ \quad c) $\frac{14}{24}$}{a) $24 : 3 = 8$, $\frac{2}{3} = \frac{16}{24}$; b) $24 : 8 = 3$, $\frac{5}{8} = \frac{15}{24}$; c) $24 : 12 = 2$, $\frac{7}{12} = \frac{14}{24}$}{}
\end{pfloesung}
\begin{pfloesung}{}
\lz{12.}{a) $30\,\%$ \quad b) $35\,\%$ \quad c) $48\,\%$}{a) $\frac{3}{10} = \frac{30}{100} = 30\,\%$; b) $\frac{7}{20} = \frac{35}{100} = 35\,\%$; c) $\frac{12}{25} = \frac{48}{100} = 48\,\%$}{}
\end{pfloesung}
\begin{pfloesung}{}
\lz{13.}{a) $\frac{6}{9}$ \quad b) $\frac{3}{12}$ \quad c) $\frac{12}{20}$}{a) $\frac{2}{3} = \frac{6}{9}$, $\frac{7}{9}$ bleibt; b) $\frac{1}{4} = \frac{3}{12}$, $\frac{5}{12}$ bleibt; c) $\frac{3}{5} = \frac{12}{20}$, $\frac{11}{20}$ bleibt}{}
\end{pfloesung}
\begin{pfloesung}{}
\lz{14.}{Nenner 12; 15; 24}{a) $\frac{1}{6} = \frac{2}{12}$ und $\frac{3}{4} = \frac{9}{12}$; b) $\frac{2}{5} = \frac{6}{15}$ und $\frac{1}{3} = \frac{5}{15}$; c) $\frac{5}{8} = \frac{15}{24}$ und $\frac{7}{12} = \frac{14}{24}$}{}
\end{pfloesung}
\begin{pfloesung}{}
\lz{15.}{Mia}{Jonas $\frac{21}{25} = \frac{84}{100} = 84\,\%$, Mia $\frac{17}{20} = \frac{85}{100} = 85\,\%$.}{}
\end{pfloesung}
\begin{pfloesung}{}
\lz{16b)}{a) $3\frac{1}{2}$ \quad b) $3\frac{2}{3}$ \quad c) $5\frac{3}{4}$}{a) $\frac{7}{2} = 3\frac{1}{2}$ ($7 : 2 = 3$ Rest 1); b) $\frac{11}{3} = 3\frac{2}{3}$ ($11 : 3 = 3$ Rest 2); c) $\frac{23}{4} = 5\frac{3}{4}$ ($23 : 4 = 5$ Rest 3)}{}
\end{pfloesung}
\begin{pfloesung}{}
\lz{17.}{a) $\frac{7}{4}$ \quad b) $\frac{17}{6}$ \quad c) $\frac{38}{9}$}{a) $\frac{1 \cdot 4 + 3}{4} = \frac{7}{4}$; b) $\frac{2 \cdot 6 + 5}{6} = \frac{17}{6}$; c) $\frac{4 \cdot 9 + 2}{9} = \frac{38}{9}$}{}
\end{pfloesung}
\begin{pfloesung}{}
\lz{18.}{a) $\frac{5}{8}$ \quad b) $2\frac{1}{4}$ \quad c) $3\frac{1}{3}$}{a) $5 : 8 = \frac{5}{8}$; b) $9 : 4 = \frac{9}{4} = 2\frac{1}{4}$; c) $20 : 6 = \frac{20}{6} = 3\frac{2}{6} = 3\frac{1}{3}$}{}
\end{pfloesung}
\begin{pfloesung}{}
\lz{19.}{$2\frac{1}{4}$\,h}{$\frac{135}{60} = 2\frac{15}{60} = 2\frac{1}{4}$\,h}{}
\end{pfloesung}
\begin{pfloesung}{}
\lz{20.}{$3\frac{1}{2}$ je Tisch}{$3\frac{1}{2}$ Baguettes: $14 : 4 = \frac{14}{4} = 3\frac{2}{4} = 3\frac{1}{2}$; jeder Tisch bekommt 3 ganze, die 2 übrigen Baguettes werden halbiert.}{}
\end{pfloesung}
\begin{pfloesung}{}
\lz{21.}{a) 15 \quad b) 5 \quad c) 7}{a) $\frac{3}{8} = \frac{15}{40}$ (mal 5); b) $\frac{28}{35} = \frac{4}{5}$ (durch 7); c) $\frac{30}{42} = \frac{5}{7}$ (durch 6)}{}
\end{pfloesung}
\begin{pfloesung}{}
\lz{22.}{a) $\frac{3}{4}$ \quad b) $\frac{3}{5}$ \quad c) $\frac{2}{3}$}{a) $\frac{24}{32} = \frac{3}{4}$ (durch 8); b) $\frac{45}{75} = \frac{3}{5}$ (durch 15); c) $\frac{36}{54} = \frac{2}{3}$ (durch 18)}{}
\end{pfloesung}
\begin{pfloesung}{}
\lz{23.}{Nenner 18; 20}{a) $\frac{5}{6} = \frac{15}{18}$ und $\frac{4}{9} = \frac{8}{18}$; b) $\frac{3}{10} = \frac{6}{20}$ und $\frac{1}{4} = \frac{5}{20}$}{}
\end{pfloesung}
\begin{pfloesung}{}
\lz{24.}{a) $4\frac{5}{6}$ \quad b) $\frac{27}{8}$ \quad c) $2\frac{3}{5}$}{a) $\frac{29}{6} = 4\frac{5}{6}$; b) $3\frac{3}{8} = \frac{3 \cdot 8 + 3}{8} = \frac{27}{8}$; c) $13 : 5 = \frac{13}{5} = 2\frac{3}{5}$}{}
\end{pfloesung}
\begin{pfloesung}{}
\lz{25.}{$\frac{7}{20} = 35\,\%$}{$\frac{14}{40} = \frac{7}{20} = \frac{35}{100} = 35\,\%$}{}
\end{pfloesung}
\begin{pfloesung}{}
\lz{26.}{nur Ole}{Nur Ole: $\frac{33}{44} = \frac{3}{4}$, genau drei Viertel. Ida: $\frac{3}{4} = \frac{27}{36}$, sie hat nur $\frac{26}{36}$.}{}
\end{pfloesung}
\end{document}
